\section*{Hoofdstuk \ref{ch:b3:glaciology} --- Glaciologie}

\begin{solution}{exo:b3:glaciology:1}
$\tau_b = \rho gH\sin\alpha = 917\times9.81\times250\times\sin\ang{3} =
\qty{118}{kPa}$, ongeveer \qty{1.2}{bar}: de grootteorde van de spanningen onder echte
gletsjers.
\end{solution}

\begin{solution}{exo:b3:glaciology:2}
(a)~$\qty{2.9e15}{m^3}\times917 = \qty{2.66e18}{kg} = \qty{2.66e6}{Gt}$;
gedeeld door \qty{360}{Gt/mm}, \qty{7390}{mm}, \qty{7.4}{m}, de vermelde waarde.
(b)~$\qty{3.0e16}{m^3}\times917/360 = \qty{76}{m}$, tegenover \qty{58}{m}.
Het ijs van Groenland rust grotendeels op vast gesteente boven het zeeniveau, dus al dat ijs zou
water aan de oceaan toevoegen. Een groot deel van het Antarctische ijs rust op ondergrond onder
het zeeniveau: dat ijs verdringt al zeewater, en het smelten ervan vervangt alleen
het water dat het verdringt; de drijvende ijsplaten voegen helemaal niets toe.
\end{solution}

\begin{solution}{exo:b3:glaciology:3}
Winst $\num{6e6}\times1.2 = \qty{7.2e6}{m^3}$, verlies $\num{4e6}\times2.1 =
\qty{8.4e6}{m^3}$ aan water: netto \qty{-1.2e6}{m^3}, een gemiddelde soortelijke balans
van $-1.2\times10^6/10^7 = \qty{-0.12}{m}$ waterequivalent. De
accumulatiezoneverhouding is $6/10 = 0.6$. De gletsjer verliest massa:
hij krimpt.
\end{solution}

\begin{solution}{exo:b3:glaciology:4}
$1 - 820/917 = \num{0.106}$ en $1 - 830/917 = \num{0.095}$: ongeveer een tiende
van het volume is nog poriënruimte wanneer de poriën afsluiten.
\end{solution}

\begin{solution}{exo:b3:glaciology:5}
$\tau_b = 917\times9.81\times300\times\sin\ang{3} = \qty{141.2}{kPa}$. Volgens
de slabstelling, $u_s = \tfrac{2A}{4}\tau_b^3H = 0.5\times\num{1.43e-16}
\times(\num{1.412e5})^3\times300 = \qty{60}{m/yr}$. Twee keer zo dik: $u_s
\propto H^{n+1} = H^4$, zestien keer zo snel. Halve helling: $u_s \propto
(\sin\alpha)^3$, acht keer zo langzaam.
\end{solution}

\begin{solution}{exo:b3:glaciology:6}
Bij \qty{-25}{\celsius} (\qty{248.15}{K}), $A = \num{3.61e-13}\exp(-60000/
(8.314\times248.15)) = \qty{8.5e-26}{Pa^{-3}.s^{-1}}$, of
\qty{2.7e-18}{Pa^{-3}.yr^{-1}}. Bij \qty{-10}{\celsius}, $\num{3.61e-13}
\exp(-27.42) = \qty{4.4e-25}{Pa^{-3}.s^{-1}} = \qty{1.4e-17}{Pa^{-3}.yr^{-1}}$.
Bij \qty{0}{\celsius} (warme tak) \qty{1.4e-16}{Pa^{-3}.yr^{-1}}. Onder
dezelfde spanningen verhouden de snelheden zich als de \omterm{def:b3:glaciology:glen}{snelheidsfactoren}:
$\num{1.43e-16}/\num{2.67e-18} = 54$.
\end{solution}

\begin{solution}{exo:b3:glaciology:7}
$h = \sqrt{2\times10^5\times10^6/(917\times9.81)} = \qty{4.7}{km}$, dicht bij
\qty{4.9}{km} van Antarctica: het plastische model geeft de dikte van een
ijskap uit alleen de breedte ervan.
\end{solution}

\begin{solution}{exo:b3:glaciology:8}
$\tau = 200/5 = \num{40}\text{ jaar}$. De lengte zou nog blijven aanpassen gedurende
meerdere decennia, een paar responstijden, nadat de opwarming stopte: de gletsjers van vandaag
hebben het klimaat van vandaag nog niet ingehaald.
\end{solution}

\begin{solution}{exo:b3:glaciology:9}
(a)~$h = \sqrt{2\times2.03\times20\times\num{8.64e6}/(917\times\num{3.34e5})}
= \qty{1.51}{m}$. Aangezien $h^2 \propto t$ duurt het \qty{3}{m} te bereiken $100\times
(3/1.51)^2 = 393$ dagen, meer dan een winter. (b)~Ijs dat veel dikker is dan
\qty{3}{m} kan niet in één winter door bevriezing groeien: het is ijs dat
meerdere winters heeft overleefd, of ijs dat door wind en stromingen in ruggen is opgestapeld.
\end{solution}

\begin{solution}{exo:b3:glaciology:10}
(a)~$H = (1025/917)\times800 = \qty{894}{m}$. (b)~Bij \qty{900}{m} is de
dikte van de \omterm{def:b3:glaciology:grounding-line}{aardingslijn} \qty{1006}{m}, $9/8$ van de vorige; de flux
stijgt met $(9/8)^5 = 1.8$. (c)~Terwijl de lijn terugtrekt over een ondergrond die
landinwaarts dieper wordt, is het ijs bij de lijn dikker, voert het meer af en dunt het ijs
erachter uit, en de lijn trekt verder terug: een positieve terugkoppeling. Een ijsplaat
die het op de ondergrond rustende ijs stut, verkleint de flux bij de \omterm{def:b3:glaciology:grounding-line}{aardingslijn} en kan
de lijn zelfs op een aflopende helling op zijn plaats houden.
\end{solution}

\begin{solution}{exo:b3:glaciology:11}
Met $h(r) = c\sqrt{R - r}$, $c = \sqrt{2\tau_0/(\rho g)}$, $V =
\int_0^R2\pi rc\sqrt{R - r}\,\mathrm dr$. Met de substitutie $s = R - r$,
\[
\int_0^R(R - s)\sqrt s\,\mathrm ds = \tfrac23R^{5/2} - \tfrac25R^{5/2} =
\tfrac4{15}R^{5/2},
\]
dus $V = \frac{8\pi}{15}cR^{5/2}$. Verdubbeling van $R$
vermenigvuldigt $V$ en het \omterm{def:b3:glaciology:sle}{zeeniveau-equivalent} met $2^{5/2} = 5.7$.
\end{solution}

\begin{solution}{exo:b3:glaciology:12}
(a)~Maart: $14.13/16.34 - 1 = \qty{-13.5}{\percent}$; september: $4.81/7.05
- 1 = \qty{-31.8}{\percent}$. (b)~In maart wordt de ijsrand grotendeels vastgehouden door
de kusten die de Noordelijke IJszee omringen, en het elke winter opnieuw
bevriezen vult het bekken weer; het dunnere ijs van de laatste jaren smelt elke zomer meer
volledig weg. En de ijs--albedoterugkoppeling werkt in de zomer: open
water absorbeert ongeveer \qty{94}{\percent} van het zonlicht, ijs weerkaatst
\qtyrange{50}{70}{\percent}, dus elk verlies van zomerijs verwarmt het water en
smelt meer.
\end{solution}

\begin{solution}{pb:b3:glaciology:1}
\textbf{1.} De kolom boven de ondergrond, met een gewicht van $\rho gH$ per eenheid van oppervlak, wordt
langs de helling naar beneden geduwd door de component $\rho gH\sin\alpha$, in evenwicht gehouden door de
schuifspanning aan de basis; voor kleine hellingen is $\sin\alpha \approx \alpha$ (de
slabstelling met de helling van het oppervlak).

\textbf{2.} $917\times9.81\times580\times0.0190 = \qty{99.1}{kPa}$;
$917\times9.81\times470\times0.0240 = \qty{101.5}{kPa}$; $917\times9.81
\times335\times0.0330 = \qty{99.4}{kPa}$.

\textbf{3.} De spanning blijft binnen \qty{2}{\percent} gelijk, terwijl de
dikte bijna halveert: de kap benadert plastisch gedrag, met $\tau_0 \approx
\qty{100}{kPa}$.

\textbf{4.} $2\times10^5/(917\times9.81) = \qty{22.23}{m}$.

\textbf{5.} $h = \sqrt{22.23\,(R - r)}$: \qty{577}{m}, \qty{472}{m} en
\qty{333}{m} op 15, 10 en \qty{5}{km} van de rand, binnen
enkele meters van de peilingen.

\textbf{6.} $\sqrt{22.23\times20000} = \qty{667}{m}$.

\textbf{7.} Zoals in \cref{exo:b3:glaciology:11}: $V = \int_0^R2\pi r
\sqrt{22.23(R - r)}\,\mathrm dr = \frac{8\pi}{15}\sqrt{2\tau_0/(\rho
g)}\,R^{5/2}$.

\textbf{8.} $\frac{8\pi}{15}\times4.715\times(\num{2e4})^{5/2} =
\qty{4.47e11}{m^3} = \qty{447}{km^3}$.

\textbf{9.} $\pi\times20^2 = \qty{1257}{km^2}$; gemiddelde dikte
$447/1257 = \qty{0.356}{km}$, \qty{356}{m}, iets meer dan de helft van de
centrale dikte.

\textbf{10.} $2^{5/2} = 5.7$ keer.

\textbf{11.} $A = \num{3.61e-13}\exp(-60000/(8.314\times263.15)) =
\qty{4.45e-25}{Pa^{-3}.s^{-1}}$, maal \qty{3.156e7}{s}:
\qty{1.40e-17}{Pa^{-3}.yr^{-1}}.

\textbf{12.} Bij \qty{10}{km}, $H = \qty{471.5}{m}$, helling $471.5/20000 =
\num{0.0236}$, $\tau_b = \qty{100}{kPa}$; $u_s = \tfrac{2A}{4}\tau_b^3H =
0.5\times\num{1.40e-17}\times10^{15}\times471.5 = \qty{3.3}{m/yr}$.

\textbf{13.} $q = \tfrac{2A}{5}\tau_b^3H^2 = 0.4\times\num{1.40e-17}
\times10^{15}\times471.5^2 = \qty{1250}{m^2/yr}$; gemiddelde snelheid $q/H =
\qty{2.6}{m/yr}$, vier vijfde van de oppervlaktesnelheid.

\textbf{14.} $2\pi rq = \pi r^2\bar b$, dus $\bar b = 2q/r = 2\times1250/10^4
= \qty{0.25}{m}$ ijs per jaar over de binnenste \qty{10}{km}.

\textbf{15.} $h = 450$ waar $R - r = 450^2/22.23 = \qty{9.11}{km}$, bij
$r = \qty{10.9}{km}$; de \omterm{def:b3:glaciology:mass-balance}{accumulatiezone} is de schijf daarbinnen, $(10.9/
20)^2 = 0.30$ van de kap.

\textbf{16.} $A$ stijgt met \num{10.2}, en ook de snelheden onder dezelfde
spanningen. Het ijs zou sneller afvloeien dan de balans aanvoert: de
kap zou dunner worden en uitspreiden totdat de spanningen, lager geworden,
opnieuw in evenwicht waren met de balans; een zachtere plastische kap heeft een kleinere vloeispanning en een vlakker
profiel.

\textbf{17.} $\num{4.47e11}\times917 = \qty{4.10e14}{kg} = \qty{410}{Gt}$.

\textbf{18.} $410/360 = \qty{1.14}{mm}$.

\textbf{19.} $356/0.5 = 711$ jaar.

\textbf{20.} Naarmate het oppervlak daalt, bevindt het zich in warmere lucht en smelt het sneller,
de hoogte--massabalansterugkoppeling; en naarmate het dunner wordt, veranderen de spanningen en de
stroming eveneens. Het uitdunnen versnelt.

\textbf{21.} $410/243 = \num{1.7}\text{ jaar}$: de hele kap is minder dan twee
jaar aan verlies van Groenland in 2010--2019.

\textbf{22.} De ijskap Arkel heeft een \omterm{def:b3:glaciology:sle}{zeeniveau-equivalent} van
\textbf{\qty[text-series-to-math=true]{1.1}{mm}} (\qty{1.14}{mm}).
\end{solution}
